% Paper encoding and complexity proof; no new Lean complexity theorem.
\documentclass[11pt]{article}
\usepackage[T1]{fontenc}
\usepackage[utf8]{inputenc}
\usepackage{lmodern}
\usepackage[a4paper,margin=25mm]{geometry}
\usepackage{amsmath,amssymb,amsthm,mathtools}
\usepackage{booktabs,longtable,array,enumitem}
\usepackage{xcolor,hyperref,xurl}
\hypersetup{colorlinks=true,linkcolor=blue!55!black,urlcolor=blue!55!black}
\setlength{\emergencystretch}{3em}
\newtheorem{definition}{Definition}[section]
\newtheorem{lemma}[definition]{Lemma}
\newtheorem{theorem}[definition]{Theorem}
\newtheorem{corollary}[definition]{Corollary}
\theoremstyle{remark}
\newtheorem{remark}[definition]{Remark}
\newcommand{\N}{\mathbb N}
\newcommand{\code}[1]{\nolinkurl{#1}}
\newcommand{\bits}{\{0,1\}^{*}}
\newcommand{\Hstar}{\mathcal H_0}
\title{Binary Encoding of the Strict-Hub Reduction\\
  \large Exact Output Length, Polynomial Construction, and Restricted Complexity}
\author{Project Research Note}
\date{7 November 2026 (Asia/Tokyo)}
\begin{document}
\maketitle

\begin{abstract}
We fix a self-delimiting binary source encoding and an explicit sparse tree
encoding for the previously verified strict-hub/three-partition semantic
correspondence. On a valid source of length $L$ the target has an exact length
formula and at most $4L+3$ bits. A total reduction on all binary strings,
including malformed inputs, is computable in $O((L+1)^2)$ time on a
deterministic multitape Turing machine. A precisely specified canonical
strict-hub language has polynomial certificates; with the classical external
strong NP-completeness theorem for 3-PARTITION, it is strongly NP-complete.
The broader tree problem inherits strong NP-hardness, but its NP membership
is not proved here. Encoding, runtime and complexity arguments in this note
are paper proofs, not Lean-checked complexity theorems.
\end{abstract}

\section{Evidence Boundary and Routing Semantics}
\label{sec:semantics}

The semantic package published in
\href{https://github.com/huam54925-cpu/uav-routing-certification/commit/dcc3db87805957680d7dbbba50f8eb6f8c7ef78e}{\code{dcc3db8}}
contains \code{HubThreePartition.exact_iff_threePartition} and the more general
\code{routing_iff_unitMatrix}. Their recorded local kernel checks and axiom
audits precede this note. No old check is relabeled as validation of the new
binary encoding or a Turing-machine model.

A route begins and ends at the sole depot, has a nonempty task sequence,
has no intermediate depot, and uses only the prescribed tree edges.
Repeated non-depot visits are permitted and every occurrence is charged.
A family of $m$ routes must realize exactly the prescribed non-depot visit
counts and directed edge multiplicities, and every route has service cost at
most $B$. There are no travel charges. This note restricts service fees and
budgets to natural numbers, interpreted as rationals by the Lean model.
It does not change the larger repository's signed-rational semantics.

\section{A Fixed Integer Code and Two Typed Formats}
\label{sec:code}

For a positive integer $t$, let $b(t)=\lfloor\log_2 t\rfloor+1$ be its bit
length and let $\operatorname{bin}(t)$ be its binary expansion without leading
zeros. Encode $n\in\N$ by gamma coding $n+1$:
\begin{equation}\label{eq:gamma}
 g(n)=0^{b(n+1)-1}\operatorname{bin}(n+1),\qquad
 \ell(n)=|g(n)|=2b(n+1)-1.
\end{equation}
Thus $g(0)=1$, $g(1)=010$, $g(2)=011$ and $g(3)=00100$.
There are no format tags, separators, alignment bits or rational-denominator
fields. The source and target are different typed formats; their binary
languages need not be disjoint.

\begin{lemma}[Unambiguous integer decoding]\label{lem:decode}
The code $g$ is prefix-free and uniquely decodable, including zero.
\end{lemma}
\begin{proof}
Count $z$ leading zeros before the first one, then read exactly $z$ further
bits. The resulting $(z+1)$-bit positive integer is $n+1$.
No first one or insufficient further bits means rejection. Codes with
different $z$ disagree at the position where one first becomes one; codes
with the same $z$ have the same length. Thus no code is a proper prefix of
another, and concatenated fields are uniquely parsed.
\end{proof}

\begin{definition}[Source format and language]\label{def:source}
A well-formed source has code
\begin{equation}\label{eq:source}
 g(m)\,g(B)\,g(a_0)\cdots g(a_{3m-1}).
\end{equation}
A valid source additionally satisfies
\begin{equation}\label{eq:valid}
 B>0,\qquad B<4a_i,\quad 2a_i<B\quad(0\le i<3m),\qquad
 \sum_i a_i=mB.
\end{equation}
The language $\mathcal P$ consists of valid sources admitting a partition
into $m$ indexed groups of exactly three items and weight $B$ each.
Malformed strings and tuples failing \eqref{eq:valid} are outside
$\mathcal P$. The $m=0$ case is allowed, with an empty item list and empty
partition; it adds only trivial instances to the usual nonempty problem.
\end{definition}

\begin{definition}[Sparse tree target format]\label{def:target}
Vertices are the integers $0,\ldots,N-1$, with depot $0$ and $N\ge1$.
For $v=1,\ldots,N-1$, specify a parent $p(v)<v$ and multiplicity $H(v)\in\N$.
Specify demands $d(v)\in\N$ and service fees $s(v)\in\N$ for all vertices,
including the depot. The code is
\begin{equation}\label{eq:target}
 \begin{aligned}
 &g(N)g(m)g(B)\,\prod_{v=1}^{N-1}\bigl(g(p(v))g(H(v))\bigr)\\
 &\hspace{22mm}\cdot\prod_{v=0}^{N-1}g(d(v))\,\prod_{v=0}^{N-1}g(s(v)).
 \end{aligned}
\end{equation}
Products denote concatenation in the displayed index order, not numerical
multiplication. No trailing bits are allowed. The strict parent inequalities
make the edges $\{p(v),v\}$ a rooted tree. Its directed prescribed flow has
multiplicity $H(v)$ on both directions of this edge and is zero elsewhere.
The entry $d(0)$ is encoded but ignored by exact-route coverage semantics;
the depot is never charged. This format represents symmetric tree-supported
flows; it is not a dense matrix or an encoding of arbitrary directed flows.
\end{definition}

The parser checks field counts against the available input length before
allocating arrays: the remaining target requires $4N-2$ integer fields after
the three headers. Each field consumes at least one bit. Huge binary values
of $N$ therefore cannot trigger an exponentially long parsing loop.
An equivalent implementation first decodes every supplied integer to EOF,
then checks that the actual payload count is $4N-2$ before using $N$ to
allocate or slice the vertex arrays. This is the strategy of the reference
codec. Parent inequalities are checked after parsing. Similarly the source parser
first reads $m,B$, decodes the remaining fields to EOF, and then compares
their actual number with $3m$; it never loops $3m$ times before this check.

\section{The Canonical Target and Exact Length}
\label{sec:length}

Given a valid source, construct $N=3m+2$, depot $0$, hub $1$, and leaf $i+2$
for source item $i$. Define
\begin{align}
 p(1)&=0,& H(1)&=m,& d(0)&=0,& s(0)&=0,\notag\\
 p(i+2)&=1,& H(i+2)&=1,& d(1)&=4m,& s(1)&=0,\label{eq:construction}\\
 &&&&d(i+2)&=1,&s(i+2)&=a_i.\notag
\end{align}
The fleet and budget headers are the original $m,B$.
The vertex renaming $0\mapsto\code{depot}$, $1\mapsto\code{hub}$,
$i+2\mapsto\code{satellite i}$ identifies this data with the canonical
instance of the existing Lean semantic theorem. This renaming and the binary
serialization themselves are specified here on paper, not as new Lean code.

\begin{theorem}[Exact size and linear bound]\label{thm:length}
Let $L$ be the length of a valid source and $L_t$ the length of its target
under \eqref{eq:source} and \eqref{eq:target}. Then
\begin{align}
 L_t&=L+\ell(3m+2)+\ell(m)+\ell(4m)+27m+4,\label{eq:exactlength}\\
 L_t&\le4L+3.\label{eq:linearlength}
\end{align}
\end{theorem}
\begin{proof}
The source length is $L=\ell(m)+\ell(B)+\sum_i\ell(a_i)$.
Since $\ell(0)=1$ and $\ell(1)=3$, count the target fields separately:
\begin{center}
\begin{tabular}{@{}ll@{}}
\toprule Fields & Length\\\midrule
Headers $N,m,B$ & $\ell(3m+2)+\ell(m)+\ell(B)$\\
Hub parent and multiplicity & $1+\ell(m)$\\
Leaf parents and multiplicities & $18m$\\
Demands, including depot & $1+\ell(4m)+9m$\\
Services, including depot and hub & $2+\sum_i\ell(a_i)$\\\bottomrule
\end{tabular}
\end{center}
Adding proves \eqref{eq:exactlength}, also at $m=0$.
The inequalities $3m+3\le4(m+1)$ and $4m+1\le4(m+1)$ imply
\[
 \ell(3m+2)\le\ell(m)+4,\qquad \ell(4m)\le\ell(m)+4.
\]
Every $a_i$ is positive, so $\sum_i\ell(a_i)\ge9m$. Consequently
\[
 \begin{aligned}
 L_t&\le L+3\ell(m)+27m+12\\
    &\le4L+12-3\ell(B)\le4L+3,
 \end{aligned}
\]
where the last inequality uses $B>0$, hence $\ell(B)\ge3$.
\end{proof}

\begin{remark}[Representation dependence]
The tree is explicit: every parent edge, multiplicity, demand and service
entry is encoded. Vertex identities are their record positions, so the
construction needs only parent labels $0$ and $1$. If instead all $N^2$
entries of a dense flow matrix are serialized, even zero entries consume
bits and the length becomes $O(L^2)$. That remains polynomial but is not
the exact formula or linear bound just proved. Adding version tags or
rational-denominator fields likewise changes the constants.
\end{remark}

\section{A Total Polynomial-Time Reduction}
\label{sec:reduction}

\begin{definition}[Target language $\Hstar$]\label{def:h0}
The language $\Hstar\subseteq\bits$ consists of well-formed target strings
with $N=3m+2$ and exactly the structural data in
\eqref{eq:construction}, except that the $3m$ leaf service fees and $B$
may be arbitrary natural numbers, for which the encoded exact-route
instance is feasible. Noncanonical or malformed targets are outside
$\Hstar$. Source strict bounds and total-weight equality are not required
in this target language. Hub-only rows, zero fees and budget zero are allowed.
\end{definition}

Let $y_{\rm no}$ be the target of the fixed source tuple
\[
 m=2,\quad B=30,\quad a=(9,9,9,13,10,10).
\]
This tuple satisfies all source hypotheses. A triple containing $13$ weighs
at least $31$; a triple not containing $13$ weighs at most $29$.
Thus it has no three-partition. The published Lean example
\code{HubThreePartitionExamples.infeasible_exact_routing} proves that its
canonical route instance is infeasible. Its source length is $54$ and
\eqref{eq:exactlength} gives $|y_{\rm no}|=129$.

Define $R:\bits\to\bits$ as follows. Parse and validate the source; on
failure return $y_{\rm no}$. On success return the encoded construction
\eqref{eq:construction}. In particular $m=0,B=0$ must take the failure
branch: the source requires $B>0$, even though $\Hstar$ admits empty fleets
at budget zero.

\begin{theorem}[Total many-one reduction]\label{thm:reduction}
For every binary string $x$,
\[
 x\in\mathcal P\quad\Longleftrightarrow\quad R(x)\in\Hstar,
 \qquad |R(x)|\le4|x|+129.
\]
There is a deterministic multitape Turing-machine implementation of $R$
running in $O((|x|+1)^2)$ time.
\end{theorem}
\begin{proof}
For valid sources, the semantic iff from Section \ref{sec:semantics} gives
the claimed equivalence. An invalid source is not in $\mathcal P$ and maps
to a canonical infeasible target, which is not in $\Hstar$.
The length bound follows from Theorem \ref{thm:length} on the valid branch
and the fixed length $129$ otherwise.

For the time bound put $L=|x|$. Decode to EOF before using the claimed
item count; the number of actual fields is at most $L$ and no decoded
integer has more than $L$ bits. Counting the fields and testing their
number against $3m$ is polynomial even if the claimed $m$ is enormous.
Partial sums have at most $L+\lceil\log_2(L+1)\rceil$ bits. The $O(L)$
schoolbook additions, shifts and comparisons therefore take $O((L+1)^2)$
bit operations. The strict bounds use constant shifts; checking total
weight uses one schoolbook multiplication $mB$, also $O((L+1)^2)$.
These are multitape operations on binary words, not unit-cost arithmetic
on unbounded integers.

On the valid branch $3m\le L$, so generating the header and hub fields
and writing $O(m)$ constant fields is polynomial. The leaf weight codewords
can be copied from the parsed source; the final output has linear length.
Sequential scans, copies and counters can all be implemented within the
same quadratic bound. Failure emits a fixed word. This proves the stated
time bound for all strings, including the empty string. A short output
alone would not have proved this algorithmic claim.
\end{proof}

\section{Polynomial Certificates for the Restricted Target}
\label{sec:np}

\begin{theorem}[Restricted NP membership]\label{thm:np}
$\Hstar$ belongs to NP under the explicit encoding above.
\end{theorem}
\begin{proof}
First parse the target and check all canonical fields, using the length
guard of Section \ref{sec:code}. Do not allocate $m$ rows from an unchecked
header. Once $N=3m+2$ and the explicit records have been verified, $m=O(L_t)$,
where $L_t$ is the input length.

The certificate assigns every leaf $i$ a route label $q_i<m$. Encode each
label in exactly
\[
 w=\max\{1,\lceil\log_2(m+1)\rceil\}
\]
bits, permitting leading zeros. The certificate has exactly $3mw$ bits,
and no trailing bits. Reject out-of-range labels. For $m=0$ it is empty.
Accept precisely when each row's sum of assigned leaf fees is at most $B$.

For soundness, visit each row's assigned leaves in any fixed order, inserting
the hub before, between and after them, and attach the original depot at
both ends. A row receiving no leaf gives the legal route $[0,1,0]$.
There are exactly $m$ routes, one visit per leaf and $m+3m=4m$ hub visits.
Every leaf edge is used once per direction and the depot edge $m$ times
per direction. Hub service is zero, so the checked row sum is the full
route cost. This gives exact feasibility.

Conversely the existing hub matrix theorem extracts a natural matrix with
each leaf column summing to one. Unit-column uniqueness gives its route
label, and the matrix's row budget is exactly the certificate's inequality.
Thus every feasible target has an accepted certificate. This uses
\code{routing_iff_unitMatrix}, which does not require source strict bounds
or total-weight equality. Empty rows must not be rejected.

The certificate length is $O(L_t\log(L_t+1))$. One implementation scans
the labels and weight records for each row, comparing binary labels and
accumulating that row's cost with schoolbook arithmetic. A conservative
bound is $O((L_t+1)^3)$ on a multitape machine, including canonical-format
validation. Too long or malformed certificates can be rejected while
reading the expected length plus one bit. Hence this is a polynomial
certificate verifier.
\end{proof}

There is also a direct size check: in any exact feasible family in $\Hstar$,
the total number of task occurrences is $4m+3m=7m$, and the total number of
vertices in the full route lists is $7m+2m=9m$. The total number of edge
occurrences is $8m$. These identities still hold with zero fees, hub-only
rows and $m=0$. They do not bound route-list length polynomially for an
arbitrary binary-encoded large-multiplicity instance outside $\Hstar$.

\section{Complexity Consequences and Their Exact Scope}
\label{sec:complexity}

The external premise is the classical strong NP-completeness of
3-PARTITION with positive integer sizes, bounds $B/4<a_i<B/2$, and total
size $mB$. Garey and Johnson give the definition and the strong
NP-completeness theorem in Section 4.2 of
\emph{Computers and Intractability}~\cite{gj79}: the 3-PARTITION definition
is on printed page 96, and Theorem 4.4 is on printed page 99.
This is an external complexity theorem, not an axiom or theorem added to
the Lean development.

The usual finite-list binary representation and \eqref{eq:source} are
polynomially interconvertible: gamma coding at most doubles each positive
integer's bit length up to a constant, and the number of explicit item
fields is part of the input size. This re-encoding preserves every numeric
value and the valid-instance conditions. Allowing the trivial $m=0$ case
does not remove hardness of the restriction $m\ge1$.

\begin{theorem}[Strong NP-completeness of the canonical sublanguage]
\label{thm:strong}
Using the external theorem just stated, $\Hstar$ is strongly NP-complete.
The general exact service-budget routing problem on trees, under this
binary sparse representation and allowing at least these nonnegative
integer inputs, is strongly NP-hard.
\end{theorem}
\begin{proof}
Theorem \ref{thm:reduction} supplies a polynomial many-one reduction from
3-PARTITION, and Theorem \ref{thm:np} gives membership of $\Hstar$ in NP.
For strong hardness one must additionally check numeric magnitudes.
On a nonempty valid source, every input size $a_i<B$ and the largest output
integer, including vertex count, demand and fleet fields, is at most
\[
 \max\{4m+2,B\}.
\]
The strongly hard 3-PARTITION subfamily has $B$ and its sizes bounded by
a polynomial in its item count $3m$. The output values are consequently
bounded by a polynomial in $N=3m+2$ as well. A fixed failure output adds
only a constant-sized exception. The reduction thus preserves a
polynomially bounded numerical subfamily, establishing strong hardness;
this step does not follow from bit-length bounds alone.

The same construction yields valid instances of the general tree problem
and retains the semantic iff, so that problem inherits strong NP-hardness.
No membership argument for that general problem is used or claimed.
\end{proof}

The hard family has tree height two and treewidth one. Its number of routes
$m$, number of vertices and hub degree grow with the input; no claim about
fixed fleet size or fixed hub degree follows. The result is a model-specific
adaptation of the standard partition construction, with the repository's
exact repeated-visit accounting. It is not a claim of a new underlying
hardness technique or a literature-priority result. Capacity-based CVRP
theorems are not silently substituted for service-per-visit semantics.

\section{Validation and Next Formalization Boundary}
\label{sec:validation}

\paragraph{Status addendum, 7 November 2026.}
The later \code{2a64246} package formalizes typed and natural-number-list
assignment checking and the independently defined numbered target's
correspondence with the existing hub model. Neither interface requires
the source \code{ValidInput} predicate. See \href{https://github.com/huam54925-cpu/uav-routing-certification/blob/2a64246acd8e5fe323e8d6bbe3fd3399cdedd3df/docs/HUB_INTERFACES_STATUS.md}{the published interface status record}.
List decoding is not binary decoding. Binary codecs, exact lengths and
machine-cost formalization remain separate; the codec/length Lean package
was explicitly deferred from that batch. Issue \#19's later A/B paper
covers arbitrary internal subtrees for supplied witnesses and unrestricted
sorted-budget matching, without a Lean interface or general realizability
claim. This status edit changes no mathematical proof.

The earlier Lean package checks the semantic theorem and finite arithmetic.
The new proofs here cover the serialization contract, total reduction,
output length, multitape construction bound, restricted certificate
verification, and complexity corollaries conditional on the cited classical
source theorem. None of these new encoding or complexity claims is asserted
as a new Lean theorem.

A local reference codec and finite tests are kept under ignored
\code{verification/runs/issue23-encoding-20261107/}. They check integer and
structured round trips, the precise length formula and its valid-input
bound, malformed/huge headers, the fixed 129-bit no-instance, and small
source/target assignment equivalence with actual route-list accounting.
They test the executable specification; they are not Turing-machine runtime
proofs, exhaustive tests of unbounded integers, or fresh Lean builds.

The final local run used Python 3.13.4 and passed 5,124 integer-code round
trips, 945 structured source/target round trips, and 943 valid-source exact
length and upper-bound checks. Small source enumeration with
$m\in\{1,2\}$ and $1\le B\le16$ covered 630 valid instances (598 feasible,
32 infeasible), 37,926 assignment comparisons and 37,929 actual list
reconstructions. All 8,191 binary strings of length at most 12 were also
checked for totalization, canonical output shape and the total length bound;
canonical shape does not assert feasibility. The run includes malformed
headers, truncated payloads, noncanonical trees, zero-budget targets,
empty rows and invalid certificates. Integer tests include selected
payloads up to 4,097 bits, not every integer of that width.
At the original \code{86993ec} publication, the reference checker was
local-only. Public source was subsequently published in \code{2a64246}.
From the repository root, the current reproduction command is:
\begin{quote}
\texttt{python tools/check\_hub\_encoding.py}
\end{quote}
The script writes to standard output. The original
\code{reference_codec_results.json} and all later generated reports stay
local under ignored \code{verification/runs/}; scope and commands are in
\code{docs/REPRODUCING_BOUNDED_CHECKS.md}.

The original Garey--Johnson book scan was fetched through the configured
local system proxy and its relevant page images were inspected directly:
printed pages 95, 96 and 99 (PDF pages 53, 54 and 55). Page 99 explicitly
preserves a numerical bound polynomial in the number of elements in its
hardness construction. The reference check therefore uses the original
source text, not just a search-result summary. The downloaded book and
page images remain local run material. Static checks of this note verify
balanced environments and braces, internal references and bibliography keys;
they do not compile or render the document.

The next formalization can separately implement an injective binary codec
with decoder round trips, prove the exact length formula, and connect an
explicit machine-cost model to the constructive bound. The external
strong NP-completeness premise must remain identified unless a compatible
formal complexity library supplies it. A general NP-membership theorem
requires its own certificate-size argument. Issue \#19 remains at its
paper contracts described in the status addendum.

This standalone English source is intended for pdfLaTeX or XeLaTeX. No
LaTeX compilation or rendered layout check is claimed in this batch.

\paragraph{Original publication scope (\code{86993ec}), 7 November 2026.}
This note and the README navigation/status update are published on the existing
\code{research/structural-service-cuts-20261023} branch after semantic commit
\code{dcc3db8}, using the configured system proxy. Publication reuses the
finite checks and source review recorded above; it adds no Lean build,
new runtime measurement, or LaTeX rendering. The local checker, generated
results, downloaded book, page images and build products are excluded.
Existing Lean source, toolchain, dependency pins and the Issue \#19
specification are unchanged. No broader issue is closed by this publication.

\begin{thebibliography}{9}
\bibitem{gj79}
M. R. Garey and D. S. Johnson.
\emph{Computers and Intractability: A Guide to the Theory of NP-Completeness}.
W. H. Freeman, 1979, Section 4.2, pp.~95--99, especially the definition
on p.~96 and Theorem 4.4 on p.~99.
\href{https://perso.limos.fr/~palafour/PAPERS/PDF/Garey-Johnson79.pdf}{University-hosted scan of the original book}.
\end{thebibliography}
\end{document}
